%% ============================================================================
%% Fichier  : annexes/Ann_E_bibliographie.tex
%% Titre    : Bibliographie Commentée
%% Version  : 0.1 -- novembre 2026
%% ============================================================================

\chapter{Bibliographie Commentée}
\label{ann:biblio}

Cette annexe dit, pour chaque référence principale, ce qu'elle apporte à
l'auditeur public et par où commencer. La liste complète figure dans les
références bibliographiques à la fin du livre.

\section*{Pour l'évaluation du contrôle interne}

\begin{itemize}
    \item \textbf{COSO, \emph{Référentiel intégré de contrôle interne}}
          (2013, traduction IFACI 2014)~\cite{coso2013icfr}. Le texte de
          base des Chapitres~\ref{chap:sci} et~\ref{chap:eval-sci}. Lire
          d'abord la synthèse, puis le chapitre sur l'efficacité du
          contrôle interne.
    \item \textbf{COSO et ACFE, guide de gestion du risque de fraude}
          (2016)~\cite{coso2016fraude}. La démarche du principe~8, utile
          pour le Chapitre~\ref{chap:donnees-fraude}.
\end{itemize}

\section*{Pour l'audit public}

\begin{itemize}
    \item \textbf{INTOSAI, GUID~5100}~\cite{intosai2019guid5100}. Les
          lignes directrices de la vérification des SI par les
          institutions supérieures de contrôle~; les treize domaines
          d'une mission.
    \item \textbf{IDI, manuel d'audit de conformité}~\cite{idi2021conformite}.
          Le plus pratique des textes de l'INTOSAI~: critères, preuve,
          documentation, rapport.
    \item \textbf{Décret n°~2013/287}~\cite{decret2013287consupe}. À
          connaître par cœur pour ses articles~2, 3, 5 et~28.
\end{itemize}

\section*{Pour l'audit interne et les référentiels}

\begin{itemize}
    \item \textbf{IIA, normes}~\cite{iia2017cripp,iia2024gias}. Les
          exigences de l'auditeur interne~; l'édition 2024 remplace celle
          de 2017.
    \item \textbf{ISACA, COBIT~2019}~\cite{isaca2018cobitintro,isaca2018cobitobjectifs}.
          À consulter comme catalogue de critères, objectif par objectif.
    \item \textbf{ISO/IEC~27001 et 27002}~\cite{iso27001_2022,iso27002_2022}.
          Les mesures de sécurité, à citer quand l'entité les a adoptées.
\end{itemize}

\section*{Pour l'analyse de données}

\begin{itemize}
    \item \textbf{Nigrini, \emph{Benford's Law}}~\cite{nigrini2012benford}.
          La référence pour appliquer la loi de Benford en audit, avec ses
          limites.
    \item \textbf{Cressey}~\cite{cressey1953}. L'origine du triangle de la
          fraude.
\end{itemize}

\section*{Pour aller plus loin}

\begin{itemize}
    \item \textbf{Minka, \emph{Vers l'auditologie}} et l'article sur la
          réponse \NEK{}~\cite{minka2026auditologie,minka2026nek}. Le
          programme de recherche dont le cours tire la cotation
          \OUI{}~/ \NON{}~/ \NEK{}~; prépublications sur HAL.
    \item \textbf{Le cours INF3226} (ENSPY), version complète de ce cours,
          pour qui veut approfondir chaque domaine.
\end{itemize}
